\documentclass[11pt]{article}

\usepackage[a4paper,margin=1in]{geometry}
\usepackage{amsmath,amssymb,amsthm}
\usepackage{mathtools}
\usepackage{enumitem}
\usepackage{microtype}
\usepackage[hidelinks]{hyperref}

% --- a bit of breathing room so the document reads well ---
\linespread{1.06}
\setlength{\parskip}{0.35em}

\theoremstyle{definition}
\newtheorem{assumption}{Assumption}
\theoremstyle{plain}
\newtheorem{proposition}{Proposition}

\newcommand{\E}{\mathbb{E}}
\newcommand{\R}{\mathbb{R}}
\newcommand{\Ug}{U_g}
\newcommand{\Um}{U_m}

\title{\textbf{Report}\\[0.4em]
\large Jean Tirole (2012), \emph{Overcoming Adverse Selection: How Public
Intervention Can Restore Market Functioning}, \emph{American Economic Review}
102(1): 29--59}
\author{Jordi Torres Vallverd\'u \\ \small MRes}
\date{\today}

\begin{document}
\maketitle

%==================================================================
\section{Introduction}
%==================================================================

Tirole asks a question that the 2007--09 crisis made important: when a market
freezes because of adverse selection, can a government jump-start it, and
if so, how? The answer is built on a single ``twist'' relative to standard
mechanism design. The textbook designer (a monopolist screening a buyer, a
regulator facing a firm) takes each agent's reservation utility as a
primitive: the agent always has a fixed outside option, and the participation
constraint is $U(\theta)\ge U_0(\theta)$ with $U_0$ exogenous. Tirole's point is
that when the designer is a government acting before a private market reopens,
this is no longer true. The outside option of a firm that declines the
government's offer is what it can get in the market, and, importantly, the market
outcome depends on who declined---that is, on the mechanism itself.
Reservation utilities are therefore endogenous to the mechanism. This
two-way feedback between mechanism design and participation constraints is the
paper's conceptual core.

The economics that follow are interesting and a bit counterintuitive. Because the
government moves first and skims off the worst (most ``toxic'') assets, it
cleans up the residual pool and so revives the very private market it is
competing against---it ``creates its own competition''. Anticipating the rebound, firms with good assets hold out for the market, so the government can never substitute fully for it: a non-comprehensive intervention with market rebound is optimal. Reviving a
frozen market is necessarily expensive (the government loses money on every type
it finances), yet---and this is the nice result---the mere existence
of the later market imposes no welfare cost. The framework rationalizes
the actual drift of TARP from asset buybacks toward equity injections, and
delivers a clean public-finance logic for when a freeze should, and should not,
be undertaken.

In this report I will try to first summarize the key fundamentals of the model, the results and the possible extensions (section~\ref{section1}).  In section~\ref{section2} I will try to connect the paper's contribution along what we have seen in the course. Finally, in section~\ref{section3} I will end with some criticisms or comments of the paper.
%==================================================================
\section{The model}
%==================================================================
\label{section1}
\subsection{Primitives}

All parties are risk neutral. A firm (seller) is cashless, protected by
limited liability, owns a legacy asset, and has a new project to
finance.

\paragraph{Legacy asset.} The asset pays $R_0$ on success and $0$ on failure.
The success probability $\theta$ is the seller's private information, drawn from
a continuous c.d.f.\ $F$ on $[0,1]$ with density $f$. Tirole assumes $F$ is
log-concave (decreasing reverse hazard rate $f(\theta)/F(\theta)$), the
standard regularity condition that keeps the screening problem well behaved.

\paragraph{New project.} The project costs $I$. If the entrepreneur
behaves it yields verifiable income $R_1$ and private benefit $b>0$; if
she misbehaves it yields no income but a large private benefit $B>b$ and
has negative social value ($B<I$). Define the social surplus of the project
\[
S \;=\; R_1 + b - I .
\]
Behaving requires leaving the entrepreneur a stake $B-b$, so the
pledgeable income is only $R_1-(B-b)$. Three assumptions organize the
analysis:

\begin{assumption}[Positive NPV]
$S>0$.
\end{assumption}

\begin{assumption}[Scope for credit rationing]
Without the legacy asset the project cannot be financed:
$R_1-(B-b)<I$, i.e.\ $S-B<0$.
\end{assumption}

\begin{assumption}[Collateral may enable financing]
For the best types collateral restores financing under symmetric information:
$R_0+(S-B)>0$.
\end{assumption}

The moral-hazard wedge $B$ is what makes the firm credit-constrained: it
cannot raise $I$ on the project alone and must sell (a stake in) the legacy
asset, whose quality only it knows. 

\paragraph{Government and welfare.} The government faces a shadow cost $\lambda$
of public funds and maximizes expected gross social surplus
\[
W \;=\; \E[U(\theta)] + \pi - (1+\lambda)D ,
\]
where $D$ is the deficit, $U(\theta)$ the seller's utility and $\pi$ the buyers'
profit. Since competitive buyers break even ($\pi=0$) in equilibrium, letting
$x(\theta)\in\{0,1\}$ indicate financing,
\[
W \;=\; (1+\lambda)\,\E[\theta R_0 + x(\theta) S] \;-\; \lambda\,\E[U(\theta)] .
\]
The first term is efficiency (assets plus financed surplus); the second is the
social cost of the rents the government must concede, weighted by
$\lambda$. Basically, a government with $\lambda>0$ is reluctant to leave large rents to
firms---this is the tension that drives the whole design problem.

\paragraph{Timing and reservation utility.} (1)~The seller learns $\theta$.
(2)~The government commits to a mechanism $C_g(\cdot)$. (3)~Each seller accepts
or declines. (4)~Competitive buyers offer menus to the remaining sellers.
(5)~Sellers choose. (6)~Payoffs realize. Laissez-faire simply deletes
stages (2)--(3). For the leading ``fleeting opportunity'' case the autarky
payoff of a seller who trades with nobody is
\[
U_0(\theta) \;=\; \theta R_0 .
\]

\subsection{The stripped-down version (buybacks only)}

To isolate the mechanism, set $R_0=1$, suppose the entire project return belongs
to the firm (no pledgeability), and take $\theta\sim U[\gamma,1]$ with
$\gamma>0$\footnote{The author will relax these assumptions and show that results mostly pass through and are general. I will discuss this in next subsection.}. A seller of type $\theta$ who sells the asset at price $p$ and
finances the project obtains $p+S$; keeping the asset (autarky) is worth
$\theta$. Hence low types sell and high types keep:
\[
\text{seller sells} \iff \theta < p+S \;=:\; \theta^\ast .
\]
Competitive buyers break even at the truncated mean of what they buy:
\[
p \;=\; \E[\theta \mid \theta<p+S]
   \;=\; \tfrac12\big(\gamma + p + S\big)
   \quad\Longrightarrow\quad
   p = \gamma+S,\qquad \theta^\ast=\gamma+2S .
\]
This is Akerlof in two lines, and it delivers the discontinuity at the heart of
a freeze:
\begin{itemize}[nosep]
  \item If $\gamma+S<I$ the break-even price is below $I$ and \textbf{there is no
  trade};
  \item If $\gamma+S\ge I$, all types $\theta<\gamma+2S$ trade at $p=\gamma+S$.
\end{itemize}
With $S<I$ and $\gamma=0$ the market fully collapses; a small piece of
good news (a rise in $\gamma$) can ``restart'' it. Conversely, a small piece of
bad news pushes $\gamma$ below $I-S$ and freezes a previously active
market---thus, there is discontinuous response to a marginal shock.

\paragraph{Adding the government.} Suppose the government posts a price $p>I$.
Were there no private market, all types $\theta<p+S$ would sell to it. But once
those types are gone, the residual pool is better, so a private market reopens at
a higher price---and anticipating this, the low types would rather
hold out. The equilibrium reconciles the two: government and market prices
are equated at $p$, types $\theta<p-S$ sell to the government, and an interval of
intermediate types trades in the rebounding private market. The government loses
money\footnote{The idea is that it overpays low types and offers them the market price to prevent them for holding up initially.}; private buyers break even. This is the ``creates its own competition''
mechanism in its simplest form.

\paragraph{When does the government intervene?} With cutoff $\theta^\ast$ and
shadow cost $\lambda$, intervening to finance the marginal type is worthwhile
only if the efficiency gain beats the extra inframarginal rents:
\[
\lambda\,F(\theta^\ast) \;<\; (1+\lambda)\,f(\theta^\ast)\,S
\quad\Longleftrightarrow\quad \lambda<1 \ \ (\text{uniform case}).
\]
So on an active market a government with $\lambda>1$ stays out. But if the
market freezes ($\gamma$ falls below $I-S$), it is optimal to let it
rebound to its minimum volume provided
\[
(1+\lambda)\,S \;>\; \lambda\,\tfrac{I+S-\gamma}{2},
\]
i.e.\ provided $I-\gamma$ is not too far above $S$. The same government switches
from laissez-faire to intervention precisely when bad news tips the
market over the ``bad cliff''.

\subsection{The general model and its main insights}
The general model relaxes uniformity ---but not log-concavity---, allows partial pledgeability (so the
designer can take a stake in the project or the asset, not just post a
price), and lets contracts be partial\footnote{Here contracts can condition payments on the legacy asset's success, i.e. specify a stake $y$ the firm retains.}. One then works with
incentive-compatible contracts and the utilities they generate: let
$\Ug(\theta)$ be type $\theta$'s rent in the government mechanism and
$\Um(\theta;\Theta_m)$ her rent in the market when the holdout set is
$\Theta_m\subseteq\Theta$. An endogenous-participation equilibrium
partitions types into those who join the government, $\Theta_g$, and the market,
$\Theta_m=\Theta\setminus\Theta_g$, with
\[
\Ug(\theta) \ge \Um(\theta;\Theta_m)\Rightarrow \theta\in\Theta_g,
\qquad
\Um(\theta;\Theta_m) > \Ug(\theta)\Rightarrow \theta\in\Theta_m .
\]
The government chooses an IC mechanism, hence a rent profile $\Ug(\cdot)$, to
maximize $W$ subject to the allocation being the (ideally unique)
equilibrium of this fixed-point problem. The fixed point is the same idea as before---$\Ug$ shapes
$\Theta_m$, which shapes the market rent $\Um$, which is the outside option in
the constraint on $\Ug$.

The optimal intervention is then characterized by six insights:

\begin{enumerate}[label=(\roman*),leftmargin=2.2em]
  \item \textbf{Pecking order.} The government buys back the weakest
  assets outright, recapitalizes (leaves on the balance sheet) intermediate
  assets, and leaves the strongest to the market. Because the government can now ask the firm to keep a stake (``skin in the game'') rather than buy outright, screening is easier and the optimal intervention is more extensive than under buybacks.
  \item \textbf{Non-comprehensive intervention / rebound.} Short of the
  prohibitively costly policy of buying everything, the market always rebounds;
  the government reduces adverse selection enough to let it restart, but
  not too much, to limit the fiscal cost.
  \item \textbf{Costly intervention.} The market already prices firms'
  willingness to fire-sell, so rejuvenation is necessarily expensive: the
  government loses money on every financed type, just like before.
  \item \textbf{No desire to shut the market down.} If the law forbids
  expropriation (firms get at least $U_0(\theta)=\theta R_0$), the voluntary-
  participation constraint can be made costless: the government would gain
  nothing from the power to ban the market, though the market deeply shapes
  how it intervenes.
  \item \textbf{When is intervention desirable?} A market failure does not by
  itself justify a bailout; the budgetary cost $\lambda$ makes the government
  reluctant. Small bad news, however, can flip the optimum from laissez-faire to
  intervention, as gains from intervening increase.
  \item \textbf{Intervention creates moral hazard.} The prospect of a
  bailout decreases ex-ante incentives to originate high-quality assets.
\end{enumerate}

A theoretical by-product (Section III) connects to the
Rothschild--Stiglitz tradition: the constrained-efficient allocation---the
one maximizing surplus subject to seller IC and buyer break-even---is an
equilibrium of the laissez-faire market game, and the unique one
surviving a ``robust choice'' refinement. The market, left alone, already
implements the best incentive-feasible allocation; the government's role is
purely to relax the adverse-selection constraint at a fiscal cost.

%==================================================================
\subsection{Extensions}
%==================================================================

Tirole proposes then several extensions; I summarize them:
\begin{itemize}[leftmargin=1.4em]
  \item \emph{Ex-ante asset quality.} Endogenizing $\theta$ via costly effort
  confirms insight (vi): the anticipated rescue lowers the marginal return to
  originating good assets, so the intervention naturally creates moral hazard. With commitment
  the optimal scheme would also internalize this effort margin.
  \item \emph{Fungibility (Myers--Majluf).} If legacy and project payoffs cannot
  be told apart, contracts reward one vs.\ two successes; the flat bottom of the
  constrained optimum (figure 4 in the paper) tilts upward but the logic is mostly unchanged.
  \item \emph{Two-sided adverse selection.} Private information on the new
  project too (i.e $S(\theta)$) makes screening two-dimensional and hard, but under perfect
  correlation it collapses to one dimension and the analysis goes through\footnote{A doubt however is what happens if $S(\theta)$ is exponentialy increasing in its argument}.
  \item \emph{A small cost of intervention} pins down a unique
  implementation: the government rescues the bottom $[\underline\theta,\theta_g]$
  and leaves $[\theta_g,\theta^\ast]$ to the market, maximizing the size of the
  free market by clearing out the worst types.
\end{itemize}

%==================================================================
\section{Placement in the course}
%==================================================================
\label{section2}

The paper sits mostly at the intersection of the four blocks we covered in class. Reading it was a good way to see how they fit together. I will try to connect them by parts:

\paragraph{Revelation principle and the informed-principal logic (Myerson 1982).}
The government's problem is a mechanism-design problem, and the revelation
principle lets us restrict attention to direct, incentive-compatible mechanisms
$(\Ug(\cdot),C_g(\cdot))$---this is what licenses working with rent profiles
rather than arbitrary contracts. What the paper changes is the
participation side. In the canonical formulation the individual-rationality
constraint is $U(\theta)\ge U_0(\theta)$ with $U_0$ fixed; here it becomes
$\Ug(\theta)\ge \Um(\theta;\Theta_m)$ with the right-hand side a
function of the mechanism. The clean separation of IC and IR that we used
throughout breaks down, and the designer's problem becomes a fixed point. This is
the single most important ``twist on the basics'' in the paper.

\paragraph{Adverse selection and market breakdown.} The
stripped-down model is the Akerlof common-values market: a seller's payoff
depends on the type of whoever ends up holding the asset, the break-even price is
a truncated mean, and incentive compatibility acts as a quantity/sorting
constraint that generates a cross-type externality---precisely the externality
that, under common values, can break decentralization (the Prescott--Townsend /
Rothschild--Stiglitz theme from the Akerlof lecture). Tirole's Section III result,
that the constrained-efficient allocation is the unique refined market
equilibrium, is the positive counterpart to that lecture's existence and
decentralization discussion.

\paragraph{Competing mechanisms / multiple principals.} The government and the
market are two principals contracting with the same population, so the paper
belongs to the competing-mechanisms world of the course (Myerson 1982 competing
hierarchies; Attar--Mariotti--Salani\'e  on non-exclusive competition for lemons).
But the placement is instructive precisely through a contrast: that
literature has the principals move simultaneously, which is exactly why
reservation utilities there are not mechanism-dependent and can be treated
as exogenous. Tirole's sequential (Stackelberg) timing is what makes the
outside option respond to the leader's mechanism. So the paper can be read as the
``competing-mechanisms'' problem with the one ingredient the standard models
assume away. The link to Attar--Mariotti--Salani\'e is quite relevant: their
result that non-exclusive competition reproduces the Akerlof outcome is what
underwrites treating the rebounding market as a competitive Akerlof market in the
first place.

\paragraph{Moral hazard.} Adverse selection is not the whole story:
the credit constraint that forces the firm to sell its asset comes from a
moral-hazard wedge $B-b$ on the new project (Assumption 2), and the
limited pledgeability $R_1-(B-b)$ is straight out of the Holmstr\"om--Tirole
``skin in the game'' logic of the moral-hazard part of the course. Insight (vi)---intervention
distorting ex-ante origination---is itself a moral-hazard result put on top
of the screening problem. The main results of the pape,  however, focus on the adverse selection problem, which is far more important in the application.
%==================================================================
\section{Critical comments}
%==================================================================
\label{section3}
In the end, I annotate some comments and doubts about the paper:
\begin{enumerate}[leftmargin=1.6em]
  \item \textbf{Adverse selection as the cause of the freeze.} The whole
  apparatus traces freezes to asymmetric information about asset quality. Tirole
  is honest (Section V) that 2007--09 also reflected fire-sale / cash-in-the-market
  pricing, funding runs, regulatory forbearance, counterparty risk and ambiguity
  aversion. A model with no leverage, no rollover risk and no interconnectedness
  cannot speak to the run dynamics that arguably mattered most for Lehman/Bear.
  The paper's value is conditional: if the binding friction is adverse
  selection, here is the optimal response. That said, the real contribution seems also
  methodological---endogenous, mechanism-dependent participation constraints---and
  this carries over to any sequential, multi-principal setting under adverse
  selection, well beyond the bailout story. That is what makes the paper matter beyond this critique.

  \item \textbf{Commitment.} The government is assumed to commit, before stage 4,
  to a scheme it will not renegotiate. Bailouts are famously time-inconsistent
  and politically renegotiated; Tirole acknowledges (Section VI) that without
  commitment take-up falls and sellers wait for better offers, as in the drawn-out
  Japanese recapitalizations. Since the no-commitment case can overturn the
  ``costless participation constraint'' result, the commitment assumption is doing
  a lot of work. Relaxing it, however, seems work for the future.

  \item \textbf{Specification of the outside option.} Several headline results,
  especially ``no desire to shut down the market,'' lean on the autarky payoff
  $U_0(\theta)=\theta R_0$ (fleeting opportunity) and on the no-expropriation legal
  constraint. Tirole argues (footnote~12) that the broad analysis still goes
  through with a less urgent need; but the specific claim that the participation
  constraint is costless leans on the fleeting-opportunity case. So the result is ensitive to how the outside option is set up.

  \item \textbf{Statics and homogeneity.} The model is static with a single
  cohort and a known type distribution $F$; there is no aggregate uncertainty
  about where $F$ sits, no learning, and no recovery dynamics. Yet the
  policy question (``is this freeze the kind that will rebound?'') is largely about
  that aggregate uncertainty.

  \item \textbf{On the positive side---and worth saying.} The pecking-order
  prediction (buy the worst, recapitalize the middle, leave the best) rationalizes
  the observed drift of TARP from outright toxic-asset purchases toward
  recapitalization (e.g.\ equity injections), and the ``government must overpay''
  conclusion gives a coherent, if politically awkward, answer to the contemporaneous
  debate (Bebchuk, Krugman, Sachs vs.\ the ``assets are undervalued, we might even
  profit'' view). The model articulates exactly why the optimistic view was wrong:
  a successful purge raises residual asset prices and makes good firms hold out, so
  the scheme must lose money. That is a counterintuitive but durable insight.
\end{enumerate}

%==================================================================
\begin{thebibliography}{9}
\small
\bibitem{akerlof1970} Akerlof, G.~A. (1970). The Market for ``Lemons'':
Quality Uncertainty and the Market Mechanism. \emph{Quarterly Journal of
Economics} 84(3): 488--500.
\bibitem{ams2011} Attar, A., Mariotti, T., and Salani\'e, F. (2011).
Nonexclusive Competition in the Market for Lemons. \emph{Econometrica} 79(6):
1869--1918.
\bibitem{coase1972} Coase, R.~H. (1972). Durability and Monopoly. \emph{Journal
of Law and Economics} 15(1): 143--149.
\bibitem{mm1984} Myers, S.~C., and Majluf, N.~S. (1984). Corporate Financing and
Investment Decisions When Firms Have Information That Investors Do Not Have.
\emph{Journal of Financial Economics} 13(2): 187--221.
\bibitem{myerson1982} Myerson, R.~B. (1982). Optimal Coordination Mechanisms in
Generalized Principal--Agent Problems. \emph{Journal of Mathematical Economics}
10(1): 67--81.
\bibitem{ps2012} Philippon, T., and Skreta, V. (2012). Optimal Interventions in
Markets with Adverse Selection. \emph{American Economic Review} 102(1): 1--28.
\bibitem{rs1976} Rothschild, M., and Stiglitz, J. (1976). Equilibrium in
Competitive Insurance Markets. \emph{Quarterly Journal of Economics} 90(4):
629--649.
\bibitem{tirole2012} Tirole, J. (2012). Overcoming Adverse Selection: How Public
Intervention Can Restore Market Functioning. \emph{American Economic Review}
102(1): 29--59.
\end{thebibliography}

\end{document}
